\section{Calibration}
\label{app:calibration}

\subsection{What is forecast}
\label{app:calib-ask}

Two probabilities are emitted and can be scored against outcomes. The ask forecast is the
posterior probability that the named target holds the named card, recorded at every ask by the
measured team. The declaration forecast is the policy's own probability that the allocation it is
about to name is correct, recorded at every voluntary declaration.

\subsection{Reliability}
\label{app:calib-decl}

See \filepath{research/v06/results/E6-calibration.txt} for the raw ten-bin reliability diagrams and
the Brier, log-loss and expected-calibration-error summaries for both forecasts.

\subsection{The calibration/resolution distinction}
\label{app:calib-versions}

\S\ref{sec:belief-exactworse} turns on a distinction that a single calibration number hides. A
forecaster can be well calibrated and poorly resolved: if it says $0.46$ and is right $46\%$ of the
time, it is calibrated, however uninformative $0.46$ is. The exact posterior and the deployed
approximation are \emph{both} well calibrated in this game --- the mean forecast of each tracks its
realised frequency closely --- and they differ in resolution, in the approximation's favour, because
the approximation carries a soft model of how opponents choose their asks and the exact object does
not. Any comparison of inference paths that reports only calibration will therefore report that they
are equally good, and be wrong.

\subsection{Caveats}
\label{app:calib-caveats}

Forecasts are collected from one team only, so the sample is not independent across seats within a
game; the intervals in the reliability table are indicative rather than inferential. The declaration
forecast is a probability under the policy's own posterior, not under the exact one, and
\S\ref{sec:belief} gives the size of the gap between them.
